\section{Introduction}

\subsection{Problem}

GraphRAG augments retrieval with a knowledge graph, but the cost structure of
graph construction limits its practicality. Table~\ref{tab:cost-compare}
contrasts a conventional \texttt{generate()}-based indexing pipeline with the
workspace-readout pipeline proposed here.

\begin{table}[t]
\centering
\caption{Indexing pipeline: conventional LLM \texttt{generate()} vs.\ J-Lens workspace readout.}
\label{tab:cost-compare}
\begin{tabular}{@{}lcc@{}}
\toprule
Dimension & Conventional (\texttt{generate()}) & This work (J-Lens readout) \\
\midrule
Time per chunk & 5--60\,s & \textbf{0.17\,s} \\
Output format & Structured JSON/triples (parse + retry) & None---direct workspace logits \\
Reliability on 7B & \texttimes\ JSON truncation (Fast-GraphRAG: ACC 3.6\%) & \checkmark\ stable readout \\
API cost & \$5--15 per 1k documents & \textbf{\$0} \\
Energy efficiency & 1$\times$ & \textbf{347$\times$} (measured, Phase 26) \\
\bottomrule
\end{tabular}
\end{table}

\subsection{Core idea}

The Jacobian Lens (J-Lens; Gurnee \& Sofroniew et al., 2026) reveals a
\emph{global workspace} inside LLMs (the \emph{J-space}): a set of verbalizable
representations that carry intermediate reasoning. The key observation behind
this work: since the model has already placed concepts and relations into its
workspace while reading a document, \emph{extracting knowledge does not require
the model to generate---it only requires reading the workspace}. Each
extraction step changes from one generation (5--20\,s) to one forward pass
(0.17\,s).

\subsection{Research questions and summary of findings}

\begin{table}[t]
\centering
\caption{Research questions and conclusions.}
\label{tab:rq}
\begin{tabular}{@{}lp{4.8cm}p{6.2cm}@{}}
\toprule
RQ & Question & Conclusion \\
\midrule
RQ1 & Can extraction steps be replaced by workspace readout? &
\textbf{Yes} (concepts, relations, roles, disambiguation), with sharp
boundaries (\S\ref{sec:negative}). \\
RQ2 & Are the extracted structures mathematically closed (grounded +
computable)? & \textbf{Yes}: grounding holds (parity with bge-m3 after
debiasing); matrix form $\equiv$ graph traversal, exactly
(\S\ref{sec:math}). \\
RQ3 & What is the computational home of each graph capability? &
One capability, one home: attributes $\to$ geometry, disambiguation $\to$
readout, multi-hop $\to$ graph/tensor (\S\ref{sec:capability}). \\
RQ4 & How much performance is retained after replacement in real frameworks? &
Both frameworks $\geq 0.9$ in both domains (LightRAG 1.21/1.23; HippoRAG2
1.25/0.95) after correcting an L4 judge-protocol artifact
(\S\ref{sec:replacement}). \\
\bottomrule
\end{tabular}
\end{table}

Table~\ref{tab:rq} summarizes the four research questions. The remainder of the
paper is organized as follows: \S\ref{sec:background} introduces the J-Lens and
the evaluation benchmark; \S\ref{sec:method} presents the extraction method,
the matrix formalization, online retrieval, and the final four-layer
architecture; \S\ref{sec:math} states the mathematical properties;
\S\ref{sec:capability} gives the capability map; \S\ref{sec:replacement}
reports end-to-end replacement experiments; \S\ref{sec:negative} collects the
negative results; \S\ref{sec:escape} describes the escape architecture;
\S\ref{sec:discussion} discusses cost accounting, methodology, limitations, and
future work.
